\documentclass[10pt,UTF8]{article}

\usepackage{geometry}

\usepackage{ctex}

\geometry{a4paper,scale=0.9}
\ctexset{today=old}

\usepackage[bookmarksnumbered=true]{hyperref}  % 生成大纲



\usepackage{times}  % 设置正文字体为 Adobe Times Roman

% \usepackage{makecell}
% \usepackage{multirow} % 表格多行合并

% \usepackage{graphicx} %插入图片的宏包
% \usepackage{subfigure}
% \usepackage{float} %设置图片浮动位置的宏包

\usepackage{amsmath}
\usepackage{amssymb}
\usepackage{mathtools}

\usepackage{physics}

% \usepackage{siunitx}
% % 关闭警告
% \ExplSyntaxOn
% \msg_redirect_name:nnn { siunitx } { physics-pkg } { none }
% \ExplSyntaxOff

\allowdisplaybreaks[3]  % 允许换页


%%%%%%%%%%%%%%%%%%%%%%%% tcolorbox %%%%%%%%%%%%%%%%%%%%%%%%%%
\usepackage[most]{tcolorbox}

\newenvironment{definition box}
{\begin{tcolorbox}[colback=blue!5!white,colframe=blue!75!black,parbox=false,before upper=\par,before lower=\par]}
{\end{tcolorbox}}

\newenvironment{theorem box}
{\begin{tcolorbox}[colback=yellow!5!white,colframe=yellow!75!black,parbox=false,before upper=\par,before lower=\par]}
{\end{tcolorbox}}
%%%%%%%%%%%%%%%%%%%%%%%%%%%%%%%%%%%%%%%%%%%%%%%%%%%%%%%%%%%%
            
%%%%%%%%%%%%%%%%%%%%% amsthm %%%%%%%%%%%%%%%%
\usepackage{amsthm}

\theoremstyle{definition}
\newtheorem{definition inner}{Definition}[section]
\newenvironment{definition}[1][]
{\begin{definition box}\begin{definition inner}[#1]\hfill\par}
{\end{definition inner}\end{definition box}}

\theoremstyle{plain}
\newtheorem{theorem inner}{Theorem}[section]
\newenvironment{theorem}[1][]
{\begin{theorem box}\begin{theorem inner}[#1]\hfill\par}
{\end{theorem inner}\end{theorem box}}

\theoremstyle{definition}
\newtheorem*{proof inner}{\textit{Proof}}
\renewenvironment{proof}[1][]
{\begin{proof inner}[#1]\hfill\par}
{\qed\end{proof inner}}

\theoremstyle{remark}
\newtheorem*{remark}{\textit{Remark}}
%%%%%%%%%%%%%%%%%%%%%%%%%%%%%%%%%%%%%%%%%%%%%%


\newcommand{\mycases}[2][0.8]{
    \left\{\vphantom{\scalebox{#1}{
        $\begin{matrix}#2\end{matrix}$
    }}\right.\!\!\!\!
    \begin{array}{l@{\quad}l}#2\end{array}
}

\newcommand{\ju}{\mathrm{j}}  % 虚数单位
\newcommand{\e}{\mathrm{e}{\mkern 1 mu}}

\newcommand{\Ev}{\,\mathcal{E}\ensuremath{v}\qty}
\newcommand{\Od}{\,\mathcal{O}\ensuremath{d}\qty}

\renewcommand{\Re}{\,\mathcal{R}\ensuremath{e}\qty}
\renewcommand{\Im}{\,\mathcal{I}\ensuremath{m}\qty}

\newcommand{\sinc}{\operatorname{sinc}} 
\newcommand{\rect}{{\mkern 1.5 mu}\operatorname{rect}\!\qty} % 矩形函数


\begin{document}

\part*{The Laplace Transform}

\section{The Laplace Transform}

\begin{definition}[The Laplace Transform]
The (bilateral) Laplace transform of a general signal $x(t)$ is defined as
\begin{equation*}
    X(s) =
    \int_{-\infty}^{\infty}x(t)\e^{-s t}\dd{t}
\end{equation*}
where $s=\sigma + \ju\omega$ is a complex variable.
\end{definition}

\begin{remark}
\textbf{Connection with the Fourier Transform}
\begin{equation*}
    \mathcal{L}\{x(t)\} = X(s) =
    X(\sigma+\ju\omega) = \mathcal{F}\{x(t)\e^{-\sigma t}\}
\end{equation*}
\end{remark}

\begin{definition}[The Inverse Laplace Transform]
The inverse Laplace transform of a function $X(s)$ is defined as
\begin{equation*}
    x(t) = \frac{1}{2\pi\ju}\int_{\sigma-\ju\infty}^{\sigma+\ju\infty}X(s)\e^{st}\dd{s}
\end{equation*}
\end{definition}
\begin{proof}
    Use the inverse Fourier transform:
    \begin{equation*}
        x(t)\e^{-\sigma t} = \mathcal{F}^{-1}\{X(\sigma+\ju\omega)\} =
        \frac{1}{2\pi}\int_{-\infty}^{+\infty}
        X(\sigma+\ju\omega)\e^{\ju\omega t}\dd{\omega}
    \end{equation*}
\end{proof}



\section{Properties of the ROC}

\begin{enumerate}
    \item The ROC consists of strips parallel to the $\ju\omega$-axis in the $s$-plane.
    \item The ROC dose not include any poles.
    \item If $x(t)$ is of finite duration and is absolutely integrable,
    then the ROC is the entire $s$-plane.
    \item If $x(t)$ is right-sided, then the ROC is of the form $\Re(s) > \sigma_0$
    for some $\sigma_0$.
    \item If $x(t)$ is left-sided, then the ROC is of the form $\Re(s) < \sigma_0$
    for some $\sigma_0$.
    \item If $x(t)$ is two-sided, then the ROC is of
    the form $\sigma_1 < \Re(s) < \sigma_2$ for some $\sigma_1$ and $\sigma_2$.
\end{enumerate}



\section{Basic Laplace Transform Pairs}

\begin{center}
\begin{tabular}
    { c @{\qquad} c @{\qquad} c}
    Signal & Laplace Transform & ROC \\
    \hline $\displaystyle
    u(t)
    $ & $\displaystyle
    \frac{1}{s}
    $ & $
    \Re{s} > 0 \vphantom{\mycases{1\\1}}
    $ \\\hline $
    -u(-t)
    $ & $\displaystyle
    \frac{1}{s}
    $ & $
    \Re{s} < 0 \vphantom{\mycases{1\\1}}
    $ \\\hline $
    \e^{-b\abs{t}}
    $ & $\displaystyle
    \frac{-2 b}{s^2-b^2}
    $ & $
    -\Re{b} < \Re{s} < \Re{b} \vphantom{\mycases{1\\1}}
    $ \\\hline $
    \e^{-a t} u(t)
    $ & $\displaystyle
    \frac{1}{s+a}
    $ & $
    \Re{s+a} > 0 \vphantom{\mycases{1\\1}}
    $ \\\hline $
    -\e^{-a t} u(-t)
    $ & $\displaystyle
    \frac{1}{s+a}
    $ & $
    \Re{s+a} < 0 \vphantom{\mycases{1\\1}}
    $ \\\hline $\displaystyle
    \frac{t^{n-1}}{(n-1)!}\e^{-a t} u(t)
    $ & $\displaystyle
    \frac{1}{(s+a)^n}
    $ & $
    \Re{s+a} > 0 \vphantom{\mycases{1\\1}}
    $ \\\hline $\displaystyle
    -\frac{t^{n-1}}{(n-1)!}\e^{-a t} u(-t)
    $ & $\displaystyle
    \frac{1}{(s+a)^n}
    $ & $
    \Re{s+a} < 0 \vphantom{\mycases{1\\1}}
    $ \\\hline $
    \delta(t)
    $ & $\displaystyle
    1
    $ & $
    \mathbb{C} \vphantom{\mycases{1\\1}}
    $ \\\hline $
    \delta(t-T)
    $ & $\displaystyle
    \e^{-s T}
    $ & $
    \mathbb{C} \vphantom{\mycases{1\\1}}
    $ \\\hline $
    t^n u(t)
    $ & $\displaystyle
    \frac{n!}{s^{n+1}}
    $ & $\displaystyle
    \Re{s} > 0 \vphantom{\mycases{1\\1}}
    $ \\\hline $\displaystyle
    \cos(\omega_0 t) u(t)
    $ & $\displaystyle
    \frac{s}{s^2+\omega_0^2}
    $ & $\displaystyle
    \Re{s} > 0 \vphantom{\mycases{1\\1}}
    $ \\\hline $\displaystyle
    \sin(\omega_0 t) u(t)
    $ & $\displaystyle
    \frac{\omega_0}{s^2+\omega_0^2}
    $ & $\displaystyle
    \Re{s} > 0 \vphantom{\mycases{1\\1}}
    $ \\\hline $\displaystyle
    \e^{-a t}\cos(\omega_0 t) u(t)
    $ & $\displaystyle
    \frac{s+a}{(s+a)^2+\omega_0^2}
    $ & $\displaystyle
    \Re{s+a} > 0 \vphantom{\mycases{1\\1}}
    $ \\\hline $\displaystyle
    \e^{-a t}\sin(\omega_0 t) u(t)
    $ & $\displaystyle
    \frac{\omega_0}{(s+a)^2+\omega_0^2}
    $ & $\displaystyle
    \Re{s+a} > 0 \vphantom{\mycases{1\\1}}
    $ \\\hline $\displaystyle
    u_n(t) = \dv[n]{\delta(t)}{t}{n}
    $ & $\displaystyle
    s^n
    $ & $\displaystyle
    \mathbb{C} \vphantom{\mycases{1\\1}}
    $ \\\hline $\displaystyle
    u_{-n}(t) = \underbrace{u(t) \ast \cdots \ast u(t)}_{n \text{ times}}
    $ & $\displaystyle
    \frac{1}{s^n}
    $ & $\displaystyle
    \Re{s} > 0 \vphantom{\mycases{1\\1}}
    $ \\\hline
\end{tabular}
\end{center}



\section{Properties of the Bilateral Laplace Transform}

\subsection{}

\begin{center}
\begin{tabular}
    { c @{\qquad} c @{\qquad} c}
    Signal & Transform & ROC \\
    \hline $\displaystyle
    x(t)
    $ & $\displaystyle
    X(s)
    $ & $
    R
    $ \\\hline $
    x_1(t)
    $ & $\displaystyle
    X_1(s)
    $ & $
    R_1
    $ \\\hline $
    x_2(t)
    $ & $\displaystyle
    X_2(s)
    $ & $
    R_2
    $ \\\hline $
    a x_1(t) + b x_2(t)
    $ & $\displaystyle
    a X_1(s) + b X_2(s)
    $ & $
    \text{At least } R_1 \cap R_2
    $ \\\hline $
    x(t-t_0)
    $ & $\displaystyle
    \e^{-s t_0} X(s)
    $ & $
    R
    $ \\\hline $\displaystyle
    \e^{s_0 t} x(t)
    $ & $\displaystyle
    X(s-s_0)
    $ & $
    R + \Re{s_0}
    $ \\\hline $\displaystyle
    x(a t)
    $ & $\displaystyle
    \frac{1}{\abs{a}} X\qty(\frac{s}{a})
    $ & $
    a R \vphantom{\mycases{1\\1}}
    $ \\\hline $
    x^*(t)
    $ & $\displaystyle
    X^*(s^*)
    $ & $
    R
    $ \\\hline $
    x_1(t) \ast x_2(t)
    $ & $\displaystyle
    X_1(s) X_2(s)
    $ & $
    \text{At least } R_1 \cap R_2
    $ \\\hline $\displaystyle
    \dv{t} x(t)
    $ & $\displaystyle
    s X(s)
    $ & $\displaystyle
    \text{At least } R \vphantom{\mycases{1\\1}}
    $ \\\hline $\displaystyle
    -t x(t)
    $ & $\displaystyle
    \dv{s} X(s)
    $ & $\displaystyle
    R \vphantom{\mycases{1\\1}}
    $ \\\hline $\displaystyle
    \int_{-\infty}^t x(\tau)\dd{\tau}
    $ & $\displaystyle
    \frac{1}{s} X(s)
    $ & $\displaystyle
    \text{At least } R \cap \{\Re{s}>0\} \vphantom{\mycases{1\\1}}
    $ \\\hline
\end{tabular}
\end{center}

\subsection{Initial- and Final-Value Theorem}

\begin{theorem}[Initial-Value and Final-Value Theorem]
    If $x(t)=0$ for $t<0$ and $x(t)$ contains no impulses or higher-order singularities at $t=0$, then
    \begin{gather*}
        x(0^+) = \lim_{s\to+\infty} s X(s) \\
        \lim_{t\to+\infty} x(t) = \lim_{s\to 0} s X(s)
    \end{gather*}
\end{theorem}



\section{The Unilateral Laplace Transform}

\subsection{}

\begin{definition}[The Unilateral Laplace Transform]
The unilateral Laplace transform of a general signal $x(t)$ is defined as
\begin{equation*}
    \mathfrak{X}(s) =
    \int_{0^-}^{\infty}x(t)\e^{-s t}\dd{t}
\end{equation*}
where $s=\sigma + \ju\omega$ is a complex variable.
\end{definition}


\subsection{Properties of the Unilateral Laplace Transform}

For causal signals, the unilateral Laplace transform is the same as the bilateral Laplace transform.

\subsubsection{}

\begin{center}
\begin{tabular}
    { c @{\qquad} c @{\qquad} c}
    Signal & Transform \\
    \hline $\displaystyle
    x(t)
    $ & $\displaystyle
    \mathfrak{X}(s)
    $ \\\hline $
    x_1(t)
    $ & $\displaystyle
    \mathfrak{X}_1(s)
    $ \\\hline $
    x_2(t)
    $ & $\displaystyle
    \mathfrak{X}_2(s)
    $ \\\hline $
    a x_1(t) + b x_2(t)
    $ & $\displaystyle
    a \mathfrak{X}_1(s) + b \mathfrak{X}_2(s)
    $ \\\hline $\displaystyle
    \e^{s_0 t} x(t)
    $ & $\displaystyle
    \mathfrak{X}(s-s_0)
    $ \\\hline $\displaystyle
    x(a t) \qquad a>0 \vphantom{\mycases{1\\1}}
    $ & $\displaystyle
    \frac{1}{a} \mathfrak{X}\qty(\frac{s}{a})
    $ \\\hline $
    x^*(t)
    $ & $\displaystyle
    \mathfrak{X}^*(s^*)
    $ \\\hline $
    \begin{matrix}
        x_1(t) \ast x_2(t) \\
        \text{both causal}
    \end{matrix}
    $ & $\displaystyle
    \mathfrak{X}_1(s) \mathfrak{X}_2(s) \vphantom{\mycases{1\\1}}
    $ \\\hline $\displaystyle
    \dv{t} x(t) \vphantom{\mycases{1\\1}}
    $ & $\displaystyle
    s \mathfrak{X}(s) - x(0^-)
    $ \\\hline $\displaystyle
    -t x(t)
    $ & $\displaystyle
    \dv{s} \mathfrak{X}(s) \vphantom{\mycases{1\\1}}
    $ \\\hline $\displaystyle
    \int_{0^-}^t x(\tau)\dd{\tau}
    $ & $\displaystyle
    \frac{1}{s} \mathfrak{X}(s) \vphantom{\mycases{1\\1}}
    $ \\\hline
\end{tabular}
\end{center}

\subsubsection{Initial- and Final-Value Theorem}
\begin{theorem}[Initial-Value and Final-Value Theorem]
    For any $x(t)$,
    \begin{gather*}
        x(0^+) = \lim_{s\to+\infty} s \mathfrak{X}(s) \\
        \lim_{t\to+\infty} x(t) = \lim_{s\to 0} s \mathfrak{X}(s)
    \end{gather*}
\end{theorem}


\end{document}